\chapter{Various Ice Phases}\label{ch:icephases}
Approximately twenty distinct three-dimensional crystalline phases exist, each characterized by specific arrangements of oxygen atoms in fixed positional relationships\cite{BartelsRausch2012}\cite{Fuentes-Landete}. The disposition of hydrogen atoms within these phases can vary, introducing the possibility of disorder. Concurrently, there are three amorphous (non-crystalline) phases, along with a multitude of computer-simulated ices that have been postulated theoretically. This profusion of ice phases can be attributed to the inherent tetrahedral configuration of water molecules within hexagonal ice, a configuration that serves as the foundation for the diverse crystal structures which emerge as this ice is subjected to increasing pressure, ultimately leading to its compression.\\
\begin{figure}[H]
\centering
\includegraphics*[width=.8\linewidth]{background/unitcell.png}
  \caption{Representative ice unit cells(generated using GenIce software)}
\label{fig:Ice unit cells}
\end{figure}
In the realm of crystalline ice phases, a consistent feature involves the intricate interplay of hydrogen bonding between water molecules, with each water molecule engaging with four neighboring counterparts. Generally, the two hydrogen atoms within each molecule exhibit a degree of equivalence, resulting in hydrogen bonds of comparable strength. This symmetrical arrangement of water molecules adheres to the principles known as the 'ice rules'\cite{OrtizAmbriz2019}. Evidently, the propensity for proton ordering—fixed positions with a consequent reduction in entropy—occurs more prominently at lower temperatures. Contrasting this, elevated pressures serve to diminish the distances separating second shell neighbors, yielding a reduced volume accompanied by heightened van der Waals effects.\\
The structural characteristics of ice phases entail an anticipated H-O-H bond angle that slightly deviates from the perfect tetrahedral angle ($109.47 \degree$), hovering around $107\degree$. Importantly, the complexities of phase transitions in multiple ice forms demand a tailored adaptation of the Clausius-Clapeyron equation, prompted by water's idiosyncratic negative expansion coefficient and its anomalous alterations in entropy with varying volume\cite{Stepanov_2007}. This equation, an essential cornerstone of thermodynamics, must be recalibrated to effectively capture the nuanced behavior exhibited by water during these intricate phase change\\
\section{IceIh}
Hexagonal ice, denoted as Ice-Ih, is a fundamental crystalline phase of ice and is considered the most common naturally occurring form of ice\cite{Baker2019}. Its name originates from the hexagonal lattice arrangement that its water molecules adopt. The structure of hexagonal ice is defined by the intricate interplay of hydrogen bonds and the unique geometry of water molecules. In hexagonal ice, each water molecule is surrounded by four neighboring water molecules, forming a tetrahedral arrangement. The oxygen atoms of water molecules form the vertices of a hexagonal lattice, while the hydrogen atoms extend outwards from the oxygen atoms. These hydrogen atoms are engaged in hydrogen bonds with oxygen atoms of adjacent molecules, resulting in a highly stable and organized structure.
The angles between the hydrogen bonds in hexagonal ice are close to the ideal tetrahedral angle of $109.47\degree$, although they are slightly compressed to about $107\degree$ due to the constraints of the hexagonal lattice. This angle compression contributes to the characteristic hexagonal symmetry of the ice\cite{Pamuk2018}.\\

\begin{figure}[H]
\centering
\includegraphics*[width=.8\linewidth]{icephases/images/iceIh.png}
\caption{Representative snapshots of IceIh}
\label{fig:IceIh}
\end{figure}


Hexagonal ice exhibits distinct layers within its structure, with water molecules forming planar layers that stack on top of each other. The arrangement of these layers is such that each layer is shifted slightly relative to the layer below, leading to an "ABABAB" stacking pattern. This stacking pattern is another hallmark of hexagonal ice's structural arrangement.
While hexagonal ice is one of the most well-studied forms of ice, it serves as a cornerstone for understanding the complexities of water's behavior in various environments. The understanding of its structural details provides insights into the broader landscape of ice phases and their role in phenomena like glacial formation, climate dynamics, and the behavior of water on other planets.

\section{IceIc}
Cubic ice, or Ice-Ic, is a metastable ice crystal formed through condensation of water vapor under specific conditions: below $-80\degree C$ at ambient pressure, in small droplets below $-38\degree C$ \cite{Murray_2006}, or through pressure reduction on high-pressure ices at 77K \cite{Ji_2011}. It often acts as a transition state between hexagonal and cubic ice due to its structural variability. The initial ice from supercooled water is stacking-disordered, making its composition complex.\\
Cubic ice's unique characteristics include slightly lower melting points and slightly higher vapor pressures compared to ice Ih. It can form rare halos around the Sun in the upper atmosphere. It's prevalent in freezing confined aqueous systems. It's preferred for small water droplets ($<15 nm$ radius) due to lower interfacial free energy. Larger cubic ice crystals slowly and irreversibly convert to hexagonal ice with heat evolution.\\
\begin{figure}[H]
\centering
\includegraphics*[width=.8\linewidth]{icephases/images/iceIc.png}
\caption{Representative snapshots of IceIc}
\label{fig:IceIc}
\end{figure}
Structurally, cubic ice has a face-centered cubic lattice with half-filled tetrahedral holes. Pure cubic ice was discovered in 2020. Like Ice-Ih, it has a moderately open low-density structure with a packing efficiency of about 1/3. Unlike Ice-Ih, its hydrogen-bonding arrangement is staggered to all neighbors. Cubic ice's layers are identical and slightly displaced, maintaining similar density to Ice-Ih.
\section{IceII}
 Ice-II emerges through the synthesis of hexagonal ice at a temperature of 198 K and a pressure of 300 MPa. Another avenue for its creation involves decompressing Ice-V at 238 K, yet it doesn't readily form through the cooling of Ice-III. This distinctive phase is known to constitute a significant proportion of icy moons like Jupiter's Ganymede \cite{Fortes_2005}.\\
The defining characteristic of Ice-II lies in its systematically ordered arrangement of hydrogen bonds. Unlike other ordered ices, Ice-II lacks a disordered counterpart. Interestingly,when heated towards its melting point, it transitions into Ice-Ih, Ice-III, Ice-V, or Ice-VI. An intriguing aspect is that the presence of minute quantities of ammonium fluoride (e.g., 0.5 mol\% $NH_4^+ F^-$) disrupts the existence of Ice-II from water's phase diagram, acting as a hydrogen-disordering agent\cite{Shephard_2018}.\\
The unit cell of Ice-II forms rhombohedral crystals, falling under the space group $R\bar{3}$. Within the crystal lattice, each water molecule establishes hydrogen bonds with four others - two as donors and two as acceptors. Ice-II can exist in a metastable state at temperatures below $\equiv$ 100 K between ambient pressure and $\equiv$ 5 GPa. However, at ambient pressure, a transformation into Ice-Ic takes place irreversibly above 160 K.
\begin{figure}[H]
\centering
\includegraphics*[width=.8\linewidth]{icephases/images/ice2_3.png}
\caption{Representative snapshots of IceII}
\label{fig:IceII}
\end{figure}

Ice-II is characterized by collapsed open hexagonal channels, a departure from Ice Ih. Its structural relationship to ice Ih can be conceptualized by detaching the columns of hexameric ice Ih rings, vertically displacing them while maintaining a right angle to their plane, rotating them around a $30 \degree$ axis, and then reconfiguring the hydrogen bonds in a more compact manner. This intricate process yields a density of $1.16 g.cm^-3$ \cite{Nakamura2015}. The hydrogen bonding within Ice-II is impeccably ordered and fixed, a feature evident in the connections surrounding the hexamers. Unlike other ordered ices like Ice-VIII, Ice-IX, Ice-XI, and Ice-XV, there is no corresponding disordered phase in Ice-II. While speculative disordered forms of Ice-II have been explored, they exhibit reduced stability due to higher Gibbs free energy and zero-point energy. Nonetheless, the possibility of partial hydrogen-bond disorder within actual Ice-II remains. This unique absence of a disordered phase can be attributed to the notable energy gap between the most and second most stable ice configurations \cite{Fan_2010}. Some hydrogen bonds within Ice-II appear bent, resulting in a weakened strength in comparison to the hydrogen bonds observed in hexagonal ice.
\section{IceIII}
Ice-III comes into existence when water is subjected to a pressure of 300 MPa and cooled to a temperature of 250 K. This distinct phase of ice exhibits stability across a relatively confined span of conditions within the P-T phase space. Figure3.1 illustrates its unit cell, forming the basis of tetragonal crystals.\\
Within the Ice-III crystal, a remarkable inter-connected-ness prevails: every water molecule engages with four others through hydrogen bonds, serving as both donors and acceptors of these interactions. A distinctive feature of Ice-III involves the presence of five-membered rings that amalgamate into bicyclo-heptamers. Its density is measured at $1.16 g.cm^{-3}$ (at 350 MPa, where water density stands at $1.13 g cm^{-3}$)\cite{Londono_1993}.

\begin{figure}[H]
\centering
\includegraphics*[width=.8\linewidth]{icephases/images/ice3.png}
\caption{Representative snapshots of IceIII}
\label{fig:IceIII}
\end{figure}

The hydrogen bonding within Ice-III mirrors that of hexagonal ice(Ice-Ih), manifesting as disordered and continually evolving. The crystal structure adopts a pseudo-cubic tetragonal framework, within which 12 water molecules come together, each playing a role in shaping the intricate structure.\\

Metastable Ice-III may be formed within the Ice-II phase space by raising the compression rate in the density-driven phase transition from hexagonal ice down to 170 K\cite{Bauer_2008}.
The architecture of Ice-III is founded on tightly woven right-handed four-fold helices, which house approximately two-thirds of the water molecules. These helices are interconnected by the remaining water molecules, which consequently experience a distinct molecular environment. This spatial arrangement creates a captivating dynamic within the Ice-III structure, underlining the intricate interplay of water molecules and their structural roles within this unique crystalline phase.

\section{IceIV}
Occasionally Ice-IV arises from the high-density amorphous ice(HDA). This generation occurs through a gradual heating process at a rate of $0.4 K \times min^{-1}$, commencing from 145 K while maintaining a constant pressure of 0.81 GPa\cite{Salzmann_2003}. It's noteworthy that a faster heating rate, such as $15 K \times min^{-1}$, predominantly yields Ice-XII. Ice-IV occupies a metastable state within the phase space encompassing Ice-III, Ice-V, and Ice-VI.
Structurally, Ice-IV adopts a rhombohedral crystal configuration, indicated by the space group 167 and Laue class symmetry -3m1. The unit cell dimensions are captured by 7.60 \AA (a, b, c; with angles measuring $70.1^\circ$, $70.1^\circ$, $70.1^\circ$) and comprise 16 molecules\cite{Engelhardt_1981}. Similar to other ice phases, every water molecule within the Ice-IV crystal engages in four hydrogen bonds, functioning as both donors and acceptors.\\

\begin{figure}[H]
\centering
\includegraphics*[width=.8\linewidth]{icephases/images/ice4.png}
\caption{Representative snapshots of IceIV}
\label{fig:IceIV}
\end{figure}

The architecture of Ice-IV emerges from a distinctive single catenated interpenetrating network, comprising hexamers that exhibit both puckered and flattish characteristics. These peculiar hexamers enable mutual penetration due to their structural design. Notably, these hexamers consist of more robustly hydrogen-bonded water molecules. The hydrogen bonds within these hexamers present varying lengths: while the penetrating hydrogen bond extends to 2.921 \AA, the water molecules involved also feature three shorter hydrogen bonds measuring 2.783 \AA. Additionally, the hydrogen bonds forming the flattish rings exhibit an extended length of 2.876 \AA. These interwoven networks aren't entirely autonomous; three-quarters of the water molecules establish a weaker hydrogen bond to the other network at 2.806 \AA, accompanied by an angle of $143^\circ$. This distinctive bond configuration leads to a dichotomy in molecular environments experienced by water molecules, delineating two unequal classes\cite{Salzmann2003_4}.\\
However, hydrogen bonding in this phase is not ordered in, in actuality, it conforms to randomness in accordance with the 'ice rules', involving two hydrogen atoms near each oxygen and one hydrogen atom on each O-O bond. This randomness contributes to a zero-point entropy proximating $3.396 J mol^{-1} K^{-1}$ \cite{Herrero_2014}.
Ice-IV stands as one of four disordered ice crystals, alongside cubic ice(Ice-Ic), Ice-XVI, and Ice-XVII, where the ordered form is intricate to produce\cite{RosuFinsen2022}. Currently, a definitive discovery of the ordered form remains elusive. The near-infrared spectrum of Ice-IV has been subjected to comparative analysis with that of hexagonal ice(Ice-Ih), yielding insights into its distinct molecular arrangement.

\section{IceV}
Ice V emerges from liquid water at 500 MPa, upon lowering the temperature to 253 K. Its unit cell resembles monoclinic crystal structure. Notably, Ice-V showcases four-, five-, six-, and eight-membered rings, forming a sub-structure with pairs of each ring type.
Clusters of seven molecules inhabit four lattice sites, with three sites containing two molecules and one site with a single molecule. These distributions confer distinct molecular environments. A chain of water molecules (O4) forms the core, constructing two condensed four-membered rings, aligned parallel to the a-axis. Parallel chains of alternating sites (O2 and O3) are connected through fewer molecules from the fourth type (O1), bridging gaps in the primary chain.\\
\begin{figure}[H]
  \centering
    \includegraphics*[width=.5\linewidth]{icephases/images/ice5.png}
    \includegraphics*[width=.3\linewidth]{icephases/images/ice5_unit_cell.png}

    \caption{Representative snapshots of IceV and a unit cell with labeled molecules}
    \label{fig:IceIh}
\end{figure}

O4 water molecules create hydrogen bonds among themselves and with O2 and O3 molecules, while they abstain from bonding with O1 molecules. Intriguingly, O1 molecules generate penetrating hydrogen bonds, traversing the eight-membered ring structures. All molecules amalgamate into a connected lattice with a density of $1.24 g .cm^{-3}$ (measured at 350 MPa, where water density reaches $1.13 g .cm^{-3}$)\cite{Hasted_1969}.
Hydrogen bonding within ice V mirrors the disorder seen in hexagonal ice, perpetually fluctuating. This characteristic yields a zero-point entropy close to $3.435 J mol^{-1} K^{-1}$\cite{Herrero_2014}. Crystal dimensions measure a = 9.22 \AA, b = 7.54 \AA, and c = 10.35 \AA, with angles at $90^\circ$, $109.2^\circ$, and $90^\circ$, respectively\cite{Kamb_1967}. This setup accommodates 28 molecules within the unit cell (at atmospheric pressure and 98 K).
Ice-V's ordered hydrogen-bonding counterpart, termed Ice-XIII (ice-thirteen), is inherently linked to its distinct structural configuration.
\section{IceVI}
Ice-VI forms from liquid water at 1.1 GPa cooled to 270 K. Its tetragonal crystal structure aligns with $P4_2/nmc$, akin to edingtonite silica. Within the crystal lattice, each water molecule engages in four hydrogen bonds, serving as two donors and two acceptors.
Ice-VI boasts tricyclo-hexamers composed of four-membered rings, contributing to its density of $1.31 g .cm^{-3}$ (measured at 0.6 GPa, where water density reaches $1.18 g .cm^{-3}$, with a minimal difference of $0.03 g.cm^{-3}$ at 81.6 $\degree C$ and 2.15 GPa).
Distinct interpenetrating networks coexist, forming a structure reminiscent of packed octahedra. Connectivity between networks occurs through equatorial water molecules in hexamers. Axial hexamer water molecules bridge hexamers while experiencing unique molecular environments, leading to corner-sharing clusters along the c direction (labeled as 'a' molecules) and hydrogen-bonded bonds in the a and b directions (referred to as 'b' molecules).\\
\begin{figure}[H]
\centering
\includegraphics*[width=.5\linewidth]{icephases/images/ice6.png}
\includegraphics*[width=.3\linewidth]{icephases/images/ice6_1.png}
\caption{Representative snapshots of IceVI and a unit cell}
\label{fig:IceVI}
\end{figure}

The hydrogen bonding within Ice-VI mirrors the disorder seen in hexagonal ice(Ice-Ih), marked by constant transformation. Two water molecule types, 'a' (20\%) and 'b' (80\%), accompanied by four hydrogen bond types (labeled 1 to 4), contribute to structural dynamics.
The tetragonal crystal measures a, b = 6.1812 \AA and c = 5.698 \AA, with angles at $90 \degree$, $90\degree$, and $90\degree$, respectively, corresponding to $D_2O$ at 1.1 GPa and 225 K, housing 10 water molecules\cite{Kuhs_1984}\cite{Chodkiewicz_2022}.
First-principles density functional theory was employed to simulate Ice-VI's vibrational spectrum, enabling a comparison with its hydrogen-ordered counterpart, Ice-XV.
\section{IceVII}
Ice-VII emerges from liquid water at pressures above 3 GPa through temperature reduction to ambient levels, following the Phase Diagram \cite{Kamb_1964}. Dynamic compression of liquid water on sub-microsecond scales also triggers its formation\cite{Myint_2018}, observed in diamond inclusions\cite{Tschauner_2018}, Earth's mantle, and solar system bodies. Laser-induced shock waves at about 5 GPa prompt immediate (around 10 nanoseconds) water-to-Ice-VII conversion \cite{Gleason_2017}.
A highly stable phase, Ice-VII's resilience spans broad temperature and pressure ranges, with relevance to giant planets and icy moons. It can be generated by decompressing ($D_2O$) Ice-VI below 95 K, under low temperature and ambient pressure conditions\cite{Klotz1999}.\\

\begin{figure}[H]
\centering
\includegraphics*[width=.8\linewidth]{icephases/images/ice7.png}
\caption{Representative snapshots of IceVII}
\label{fig:IceVII}
\end{figure}

The Ice-VII unit cell adopts a cubic crystal structure ($Pn\bar{3}m$, 224; Laue class symmetry m-3m), featuring two interpenetrating cubic ice lattices. Hydrogen bonds thread through water hexamers, with no interconnecting bonds, yielding a density of about $1.65 g.cm^{-3}$ (measured at 2.5 GPa and 25 $\degree C$)\cite{Hasted_1969}.
Cell dimensions of 3.3501 \AA at 2.6 GPa and 22 $\degree C$ accommodate two water molecules. Hydrogen bonding mirrors hexagonal ice's disorder, but Ice-VII undergoes a proton disorder-order transition to Ice-VIII around 5 $\degree C$, while maintaining metastability down to 77 K\cite{Jorgensen_1985}.
Theoretical predictions highlight plastic water ices between Ice-VII and liquid water, where $H_2O$ molecules rotate freely around fixed lattice sites\cite{Hernandez2018}. Though not experimentally verified, these insights expand our understanding of intricate ice behaviors.
\section{IceVIII}
Ice-VIII emerges from Ice-VII upon temperature reduction. Ice-VII transitions to Ice-VIII through a proton disorder-order shift when cooled around 5 $\degree C$, with both sharing near-identical structures aside from proton ordering. This transition leads to a slight distortion in the Ice-VII cubic lattice, causing a shift to a tetragonal crystal structure ($I4_1/amd$, 141; Laue class symmetry 4/mmm). Here, all water molecules form hydrogen bonds with four neighbors, two as donors and two as acceptors.

\begin{figure}[H]
\centering
\includegraphics*[width=.8\linewidth]{icephases/images/ice8.png}
\caption{Representative snapshots of IceVIII}
\label{fig:IceVIII}
\end{figure}
Analogous to Ice-VII, Ice-VIII features two interpenetrating cubic ice lattices, with each oxygen atom encompassed by its eight nearest neighbors. However, these atoms are connected solely to four tetrahedrally coordinated neighbors through hydrogen bonds. Notably, Ice-VIII's density is approximately $1.66 g .cm^{-3}$ (at 8.2 GPa and 223 K)\cite{Hasted_1969}, marginally less than twice cubic ice density due to longer intra-network O-O distances, accommodating interpenetration. Its molar volume is slightly smaller (by $0.65 mm^3 mol^{-1}$) than Ice-VII along the phase transition line. Triple points exist between Ice-VI and Ice-VII (5 $\degree C$, 2.1 GPa) and Ice-VII and Ice-X (100 K, 62 GPa)\cite{Fukui2022}.\\
With cell dimensions of a, b = 4.4493 \AA, c = 6.413 \AA (90\degree, 90\degree, 90\degree; $D_2O$, at 2.6 GPa and 22 $\degree C$), each unit cell accommodates eight water molecules. Uniform molecular environments prevail among all molecules. The H-O-H angle's limited variance from isolated molecules results in non-straight hydrogen bonds\cite{Besson_1994}.
\section{IceIX}
Ice-IX represents the equilibrium structure of Ice-III at lower temperatures, slightly denser and arranged within the $P4_12_12$ space group\cite{Mirzaei_2019}. At 165 K and 280 MPa, its cell dimensions are 6.692 \AA (a) and 6.715 \AA (c) . Ice-III, featuring a hydrogen-disordered structure, thrives between -24 $\degree C$ (249 K) and -44 $\degree C$ (229 K) and 240 to 340 MPa\cite{Londono_1993}. In 1967, the gradual transformation of Ice-III into the ordered hydrogen-bonded structure of Ice-IX was observed between -65 $\degree C$ (208 K) and -108 \degree C (165 K) within this pressure range\cite{Mirzaei_2019}. Notably, Ice-II irreversibly forms at intermediary temperatures.\\
Upon hydrogen ordering, the 'a' lattice constant expands while the 'c' lattice constant contracts. This results in negative thermal expansion due to torsional strain accumulation within the a-b plane and the expansion of the helical "springs" within the structure upon heating\cite{Sharif2021}.

\begin{figure}[H]
\centering
\includegraphics*[width=.8\linewidth]{icephases/images/ice9.png}
\caption{Representative snapshots of IceIX}
\label{fig:IceIX}
\end{figure}

Ice-IX is metastable within the Ice-II phase space and readily converts to Ice-II upon warming, rather than reverting to Ice-III. The shift from proton disordered starts within Ice-III and completes at lower temperatures, marked by a first-order transition near 126 K\cite{Knight2006}. The hydrogen-bonding in Ice-III largely transitions to proton-ordered in the process of becoming Ice-IX, achieved through rapid cooling in liquid nitrogen (77 K) to evade Ice-II formation. Ice-IX and Ice-III share identical structures and densities, differing only in proton ordering\cite{LaPlaca1973}. Employing density functional theory-based methods, the vibrational spectrum of Ice-IX has been explored . Based on second-order Møller–Plesset perturbation (MP2) theory, the phase transition boundary between Ice-IX and Ice-XIII has been predicted at $\equiv 0.30 GPa$ and 154 K\cite{Xu_2020}.

\section{IceX}
As pressure increases and the O-O distance contracts, Ice-VII undergoes a continuous transition to cubic Ice-X. This transition maintains the $Pn\bar{3}m$ space group, with a, b, c = 2.78 \AA at 62 GPa and 300 K\cite{Hemley_1987}. In Ice-X, the ice protons are evenly spaced and bonded between oxygen atoms within a molecular crystal\cite{Hirsch_1984}. Each oxygen atom is surrounded by four equidistant tetrahedrally placed nearest neighbor hydrogen atoms. Oxygen atoms adopt a body-centered cubic arrangement, each with eight nearest neighbor oxygen atoms forming a cube.\\
Hydrogen atoms, on the other hand, possess 12 nearest neighbor hydrogen atoms arranged in a body-centered truncated cube pattern. Ice-X is likely to be partially ionized, with both hydrogen and oxygen atoms fully ionized. Additionally, some protons are expected to occupy octahedral cavities, a behavior observed in Ice-VII at lower pressures.Ice-X looks identical as seen from the x, y or z-direction\cite{Guthrie_2013}.

\section{IceXI}

Ice-XI is the equilibrium structure of hexagonal ice at low temperatures, achieved by carefully preparing dilute K-OH solutions and cooling them just below 72 K (at ambient pressure; $D_2O$ 76 K) for approximately a week. This transition is facilitated by localized proton hopping, which becomes more effective with increased pressure\cite{Castro_Neto_2006}. The presence of hydroxide ions from K-OH creates defects in the hexagonal ice, allowing protons to move more freely between oxygen atoms and breaking the regular hydrogen bonding pattern.\cite{Fukazawa1998}//
Ice-XI exhibits ordered hydrogen bonds due to these effects and breaks the conventional hydrogen bonding arrangement seen in Ice-Ih. This ordered structure results in reduced entropy and increased stability. The $K^+$ ions from the K-OH solution occupy interstitial sites within the hexagonal ice lattice, further influencing its properties.
\begin{figure}[H]
\centering
\includegraphics*[width=.8\linewidth]{icephases/images/ice11.png}
\caption{Representative snapshots of IceXI}
\label{fig:IceXI}
\end{figure}
Ice-XI takes on the form of proton-ordered hexagonal ice, forming orthorhombic crystals with a unit cell dimension of a = 4.5019 \AA, b = 7.7978 \AA, c = 7.3280 \AA (90\degree, 90\degree, 90\degree, 8 molecules)\cite{Leadbetter_1985}. It has been found in ancient Antarctic ice and pure water samples. The transformation temperature from Ice-Ih to Ice-XI has been determined to occur at 58.9 K, while the reverse transformation from Ice-XI to Ice-Ih occurs at 73.4 K. Interestingly, Ice-Ih transformed to Ice-XI and back retains hydrogen-ordered domains and can more easily transition between the two structures. This transformation behavior has implications for ferroelectric structuring within crystalline ice films, allowing for stable proton ordering up to around 175 K, suggesting its widespread presence in various environments\cite{Sugimoto2016}.

\section{IceXII}
Ice-XII is an intriguing structure that can be obtained by heating high-density amorphous ice under certain conditions. This transformation occurs at a constant pressure of 0.81 GPa, with the temperature raised from 77 K to approximately 183 K at a rate of $\ge 15 K min^{-1}$\cite{Salzmann_2003}. The resulting Ice-XII can then be recovered at atmospheric pressure and 77 K [2096]. It's worth noting that a slower heating rate ($\ge 0.4 K min^{-1}$) tends to preferentially produce Ice-IV. Ice-XII exists in the metastable region between the Ice-V and Ice-VI phases and takes on a tetragonal crystal structure.\\
Within the crystal lattice of Ice-XII, each water molecule forms hydrogen bonds with four neighboring water molecules – two as donors and two as acceptors. What distinguishes Ice-XII is its unique screw-type hydrogen-bonded arrangement, forming a right-handed double helix. This arrangement is unlike those observed in other forms of crystalline ice. Although previous reports indicated a ring size of five molecules\cite{Lobban_1998}, recent findings suggest the presence of seven-membered rings\cite{O_Keeffe_1998}. The crystal structure contains 12 water molecules, with one-third of them experiencing a more regularly tetrahedral configuration, leading to distinct molecular environments.

\begin{figure}[H]
\centering
\includegraphics*[width=.8\linewidth]{icephases/images/ice12.png}
\caption{Representative snapshots of IceXII}
\label{fig:IceXII}
\end{figure}

Ice-XII has a density of $1.30 g .cm^{-3}$ at 127 K and ambient pressure, slightly higher than that of Ice-V $1.23 g .cm^{-3}$. Similar to other ice structures, the hydrogen bonding in Ice-XII is disordered and continuously changing. Its tetragonal crystal unit cell has dimensions a, b = 8.276 \AA and c = 4.027 \AA (90\degree, 90\degree, 90\degree)\cite{Koza_1999}. Ice-XII is metastable within the Ice-V phase space\cite{Tonauer2021}. The ordered hydrogen-bonding structure of Ice-XII is known as Ice-XIV.

\section{IceXIII}
Ice-XIII is a notable proton-ordered variant of Ice-V, characterized by its distinctive structure. The formation of Ice-XIII involves doping the ice with 10 mM HCl (approximately one molecule for every 5000 water molecules) at temperatures below 130 K and under a pressure of 500 MPa, facilitating the phase transition\cite{Salzmann_2006}. Alternative dopants have also been explored in attempts to induce this transition\cite{RosuFinsen2018}. Cycling between temperatures of 80 K and 120 K demonstrates a reversible transformation between Ice-XIII and its metastable counterpart, Ice-V\cite{Salzmann_2006_13}. The structure of Ice-XIII, specifically monoclinic $P2_1/a$ crystals, reveals cell dimensions of a = 9.24 \AA, b = 7.47 \AA, and c = 10.30 \AA (90\degree, 109.7\degree, 90\degree; $D_2O$, at atmospheric pressure and 80 K). Each unit cell encompasses 28 water molecules, distributed among seven distinct types, with four molecules representing each type\cite{Salzmann_2006}.
\begin{figure}[H]
\centering
\includegraphics*[width=.8\linewidth]{icephases/images/ice13.png}
\caption{Representative snapshots of IceXIII}
\label{fig:Ice-13}
\end{figure}
The influence of doping with various substances such as HCl, HF, LiOH, and KOH on this phase transition has been extensively investigated through crystallographic studies\cite{Salzmann_2021}. These experiments help provide insights into the underlying mechanisms of the transformation process. The phase transition boundary between ice phases Ice-IX and Ice-XIII has been theoretically predicted using the second-order Møller–Plesset perturbation (MP2) theory. This transition is anticipated to occur around 0.30 GPa and 154 K\cite{Xu_2020}.
\section{IceXIV}:
Ice XIV emerges as the proton-ordered counterpart of Ice-XII, facilitated by the incorporation of 10 mM HCl (roughly one molecule per 5000 water molecules, or DCl with D2O) at temperatures below 118 K and pressures of 1.2 GPa, serving as catalysts for the phase transition\cite{Salzmann_2006}. While this transformation remains reversible upon temperature adjustments, it is noteworthy that reduced proton-ordering is evident at lower pressures \cite{Salzmann_2006_12}. The resulting crystals adopt an orthorhombic $P2_12_12_1$ symmetry, featuring cell dimensions of a = 8.350 \AA, b = 8.139 \AA, and c = 4.083 \AA(90\degree, 90\degree, 90\degree; D2O, at atmospheric pressure and 80 K). Each unit cell comprises 12 water molecules, where 8 of the 24 proton positions remain predominantly disordered\cite{Salzmann_2006}.\\

\section{IceXV}
Ice-XV represents a hydrogen-bonding ordered variation of the Ice-VI phase, characterized by a subtle expansion. The conceptualization of its existence stems from the presence of ordered configurations observed in other low-temperature ices. The comprehensive elucidation of its structure, particularly in its fully deuterated form ($D_2O$), was established in 2009\cite{Salzmann_2009}. The realization was facilitated through the application of 10 mM DCl (deuterated hydrochloric acid, approximately one molecule per 5000 water molecules) as a dopant, promoting the rearrangement of hydrogen bonding at sub-130 K temperatures. Notably, alternative dopants have undergone experimentation as well\cite{RosuFinsen2018}. In $H_2O$ Ice-VI, hydrogen ordering proceeds more swiftly via HCl than in $D_2O$ Ice-VI, potentially involving ferroelectric intermediates in the transition process \cite{Shephard2015}.\\
Distinguished by its antiferroelectric nature, Ice-XV showcases water molecule dipoles that effectively cancel out due to the opposing alignment of hydrogen bonding within the two interpenetrating networks. While the potential for ordered yet ferroelectric hydrogen bonding akin to Ice-VIII exists, the identification of such crystals remains outstanding. This intricate interplay of ordering and alignment adds to the diverse tapestry of ice phases and their characteristics, contributing to our understanding of water's behavior in distinct conditions.\\
When subjected to elongation or negative pressure, ice can give rise to (meta-)stable ultra-low-density forms. These ices exhibit the potential for stability even at ambient pressure when accompanied by interstitial molecules that counteract structural collapse. However, within the realm of negative pressure—specifically, under metastable conditions concerning a mixed gaseous water/hexagonal ice environment—the necessity for these stabilizing molecules might be obviated. Currently, this intriguing category comprises three distinct variants. Among them, two, namely Ice-XVI and Ice-XVII, have undergone empirical verification, while the third, an empty clathrate labeled s-III, is well-supported through computational modeling yet remains pending experimental validation.

\section{IceXVI}
Initially proposed through theoretical considerations \cite{Kosyakov2001}, Ice-XVI emerges through the treatment of neon hydrate, a clathrate with CS-II structure. The generation of Ice-XVI involves an intricate process wherein continuous vacuum pumping is administered for five days, effectuating the complete removal of neon atoms that had previously occupied the structure \cite{Falenty_2014}. This transformative endeavor causes the hydrogen bonds to extend, triggering a consequential expansion of Ice-XVI. The amplification is characterized by the 3.9\% increase in the $5^12$ cavities and the 3.3\% expansion in the larger $5^{12}6^4$ cavities, a phenomenon driven by the dissipation of attractive van der Waals interactions. Among the ensemble of experimentally recognized crystalline water phases, Ice-XVI claims the distinction of being the least dense, registering a density of $0.81 g.cm^{-3}$. While anticipated to assume stability as the low-temperature phase of water under negative pressure conditions, it regrettably succumbs to collapse and decomposition at 145 K. An intriguing correlation is observed in the increment of its lattice constants, marking a minute 0.1\% escalation in comparison to the filled hydrate at low temperatures\cite{Falenty_2014}.

\begin{figure}[H]
\centering
\includegraphics*[width=.8\linewidth]{icephases/images/ice16.png}
\caption{Representative snapshots of IceXVI}
\label{fig:Ice16}
\end{figure}

The structural blueprint of Ice-XVI harmoniously mirrors the cage configuration found within the CS-II clathrate hydrate, an illustration depicted on the left (omitting water hydrogen atoms). The cubic crystals house a total of sixteen $5^{12}$ cavities, along with an additional eight more substantial $5^{12}6^4$ cavities, encompassing 136 $H_2O$ molecules per unit cell. Of note, the tetrahedral $5^{12}6^4$ cavities convene in an open tetrahedral network, with their placements evoking recollections of the cubic ice structure, interspersed by groupings of three $5^{12}$ cavities.\\
Intriguingly, Ice-XVI departs from the norm seen in filled CS-II hydrates, as it exhibits a phenomenon of negative thermal expansivity below 55 K, akin to certain other ice formations, such as hexagonal ice\cite{Yagasaki_2016}. Amid these idiosyncratic attributes, the hydrogen bonding within Ice-XVI adheres to the 'ice rules', embracing a random arrangement characteristic of these crystalline water phases.

\section{IceVII}
Molecular hydrogen ($H_2$) readily forms diverse clathrates, wherein $H_2$ molecules adeptly occupy various cages, including multiple $H_2$ within the larger ones. Under high pressures (around 400 MPa, 280 K), a distinct structure emerges—a filled intercalated ice, rather than the conventional clathrate. This structure, denoted as $(H_2O)_2H_2$, features interpenetrating spiral chains of water molecules, bearing a topological resemblance to the mineral quartz\cite{Strobel_2016}. Subsequent heating under vacuum at temperatures of 110-120 K for about an hour releases all the enclosed hydrogen, yielding a novel metastable hexagonal ice configuration. This newly formed ice adopts space group $P6_122$, with cell dimensions of 6.3305 \AA, 6.3305 \AA, and 6.0580 \AA (a, b, c; 90\degree, 90\degree, 120\degree; 6 molecules; 75 K). Its hydrogen bonding follows the 'ice rules', characterized by randomness. Above 130 K, the sample undergoes a transformation to Ice-Ih. Notably, this ice exhibits microporosity due to its hexagonal helical pores, featuring a channel diameter of approximately 6.10 \AA from the center of the oxygen atoms and a minimum diameter of around 3 \AA. If exposed to molecular H2 (about 1.8 \AA in diameter), the ice will re-adsorb it. The structure of the heavy water variant of this ice has been elucidated through neutron diffraction techniques\cite{del_Rosso_2016}.
\begin{figure}[H]
\centering
\includegraphics*[width=.8\linewidth]{icephases/images/ice17.png}
\caption{Representative snapshots of Ice-VII}
\label{fig:Ice17}
\end{figure}

\section{IceVIII, IceIX and IceXX}

Recently, experimental achievement of ice states at exceedingly high pressures beyond Ice-X has unveiled an array of possibilities, albeit with some modeling uncertainty. Given that modeling of lower pressure ices doesn't match well with experimental structural data, these outcomes are, at best, suggestive rather than definitive. A novel superionic phase(Ice-XVIII) has been proposed, approximately characterized by a triple point at 1000 K and 40 GPa, involving liquid (supercritical and ionized) water alongside high-temperature Ice-VII (around 1500 K) \cite{Goncharov_2009}. In this scenario, hydrogen ions (protons) were anticipated to exhibit remarkable mobility akin to a liquid, coursing through the solid oxygen ion lattice. A strikingly high ionic conductivity exceeding $100 S \times cm^{-1}$, comparable to metallic conductivity, is predicted.
Another superionic ice phase has also been postulated, Ice-XX. Employing a blend of synchrotron X-ray diffraction and optical spectroscopy, investigations of $H_2O$ within a laser-heated diamond anvil cell up to 150 GPa and 6500 K have uncovered intriguing transitions. Above 900 K and 20 GPa, a first-order transition leads to ice characterized by body-centered cubic (bcc) oxygen atoms, coexisting with a superionic hydrogen fluid. Subsequently, above 1300 K and 29 GPa, fcc superionic ice-eighteen emerges. Its density, while within its stability region ($\equiv 30 GPa, 1250 K$), amounts to $1250 kg \times m^{-3}$.
\section{Clathrates, LDA and HDA}
Gas hydrates are crystalline solids with ice-like properties made
of water (the host) and gas (the guest) molecules. Typically,
gas hydrates form at low temperatures (about $0 \degree C$) and
moderate pressures (between 1 and 10 MPa) \cite{SloanJr2007} . One crucial
distinction is that gas-hydrates are not chemical compounds
and do not have strong chemical bonds between water
molecules and gas molecules. In 1948, Powell identified gas
hydrates as "Clathrate chemicals" The Latin term "Clathratus",
which means "to en-cage", is the source of the English word "clathrate". However, "Gas hydrates" \cite{Ke2019} or "Hydrates" are most frequently used for simplicity's sake. 
The van der Waals interaction between "guest" and "host" molecules is what causes gas hydrate formations to stabilize. They can be divided into three groups: structure-I (sI), structure-II (sII), and structure-H (sH), which are typically determined by size, the form of "guest" molecules, and favorable thermodynamic conditions(low temperature, high pressure, etc.) \cite{Lee2019}\cite{Khan2019} . It should be noted that both sI and sII contain cages (referred to as $5^{12}$ or "little cages") made out of 12 planar pentagonal faces.Furthermore, sI has $5^{12} 6^2$ (big cage), but sII has $5^{12}6^4$ (large cage) \cite{Walsh2009} . Methane ($CH_4$ ) and ethane ($C_2H_6$ ) create sI-type hydrates, whereas propane ($C_3 H_8$ ) and iso-butane, which are somewhat bigger molecules, create sII-type hydrates. The guest molecules neohexane and methylcyclopentane can be contained within the hexagonal structure of sH, which is unlike the cubic structures shared by sI and sII and that of sH \cite{Lee2019}.
\begin{figure}[H]
\centering
\includegraphics*[width=.8\linewidth]{icephases/images/cages.png}
  \caption{(a) $5^{12}$, $6^25^{12}$ and $6^45^{12}$ cages in stick representation, (b) an instance of blob like clathrate formation using those cages as backbone}
\label{fig:Ice-13}
\end{figure}
Over millions of years, deep-ocean (marine) and permafrost environments generate deposits of highly concentrated gas-hydrate. They are viewed as highly promising fuel resources with a wide range of scientific and industrial uses, including the modeling of climate change, the storage of carbon dioxide, the extraction of hydrocarbons, the storage of hydrogen and natural gas, the development of separation and
refrigeration technologies, the study of marine biology, and the planetary surface chemistry \cite{SloanJr2007}\cite{Walsh2009}\cite{Boswell2011}\cite{Bhattacharjee2019}\cite{Bhattacharjee2016}.\\
Amorphous ice exists in three distinct forms: Low Density Amorphous Ice (LDA), Medium Density Amorphous Ice(MDA) and High Density Amorphous Ice (HDA)\cite{Silonov2015}.
LDA is a state of amorphous ice that exhibits a lower density compared to its crystalline counterpart, Ice-Ih. It is formed by rapidly cooling water vapor below its freezing point. In LDA, the water molecules lack long-range order, resulting in a disordered, glass-like structure. Despite its lack of a well-defined lattice, LDA retains a degree of tetrahedral coordination typical of water molecules. This phase is notable for its intriguing properties, including enhanced molecular mobility and a higher propensity for molecular reorientation compared to other amorphous ice forms. LDA is often considered a prototypical example of the dynamic behavior exhibited by amorphous materials.
\begin{figure}[H]
  \centering
    \includegraphics*[width=.8\linewidth]{icephases/images/clathrates.png}
    \caption{Representative snapshots of (a)sI, (b)sII and (c)sH crystal.}
    \label{fig:Ice-13}
\end{figure}
Conversely, HDA is characterized by a higher density and is produced by subjecting high-pressure amorphous ice to lower temperatures. This results in the compression of water molecules into a more tightly packed arrangement. HDA's structure lacks long-range order and possesses a denser and more compact packing compared to LDA. It exhibits properties distinct from both LDA and crystalline ice, such as its ability to relax into LDA upon heating, known as the LDA-HDA transition. HDA's denser structure has been associated with a reduced molecular mobility, contributing to its higher stability at lower temperatures.

% \begin{figure}[H]
% \centering
% \includegraphics*[viewport= 80 130 725 505, width=.4\linewidth]{usingLatex/images/MobileMarketShare2006.pdf}
% \caption{Mobile Phone Suppliers, Market Share 2006 (Millions of
% Units Shipped)}
% \label{fig:using:MobilePhoneSuppliers2006MarketShare}
% \end{figure}


% \begin{figure}[H]
% \begin{minipage}[t]{7.2cm}
% \begin{center}
% \includegraphics*[viewport= 80 130 725 505, width=.8\linewidth]{usingLatex/images/MobileMarketShare2006.pdf}
% \caption{Caption for Figure}
% \label{fig:using:Example1}
% \end{center}
% \end{minipage}
% \hfill
% \begin{minipage}[t]{7.2cm}
% \begin{center}
% \includegraphics*[viewport= 80 130 725 505, width=.8\linewidth]{usingLatex/images/MobileMarketShare2006.pdf}
% \caption{Caption for Figure}
% \label{fig:using:Example2}
% \end{center}
% \end{minipage}
% \end{figure}


% \begin{figure}[H]
% \begin{center}
% \includegraphics[width=.34\linewidth]{usingLatex/images/samplepng}
% \caption{Caption for Bitmap Image Example} \label{fig:using:samplepngImage}
% \end{center}
% \end{figure}
%

